% Part: proof-theory
% Chapter: sequent-calculus
% Section: admissible-derivable

\documentclass[../../../include/open-logic-section]{subfiles}

\begin{document}

\olfileid{pt}{seq}{adm}

\olsection{Aturan Admisibel lan Aturan sing Bisa Diturunake}

Akeh asil teori bukti ngenani---utawa nggunakake---asil bab apa
kabeh sing bisa !!{prove} nganggo aturan tartamtu uga bisa
!!{prove} tanpa aturan kasebut. Teorema eliminasi potongan minangka
conto sing paling kondhang. Teorema iki nuduhake yen saben sekuen sing
bisa !!{prove} nganggo aturan sistem sekuen kanthi~\Cut{} uga bisa
!!{prove} tanpa~\Cut. Amarga sipat iki wigati banget, sipat kasebut
diwenehi jeneng.

\begin{defn}
Aturan~$R$ diarani \emph{admisibel} ing sistem sekuen yen saben sekuen
sing bisa dibuktekake ing sistem kasebut kanthi tambahan aturan~$R$
uga bisa dibuktekake tanpa~$R$.
\end{defn}

\begin{ex}\ollabel{ex:land-deriv}
Aturan $\LeftR{\land}$ duwe wujud beda ing \Log{G1c} lan
\Log{G3c}. Ing \Log{G1c} ana rong versi aturan iki: siji kanthi
konjung kiwa $!A$ saka !!{formula} utama~$!A \land !B$ ing
premis, lan sijine kanthi konjung tengen~$!B$:
\[
\Axiom$ !A, \Gamma \fCenter \Delta$
\RightLabel{$\LeftR{\land}_1$}
\UnaryInf$ !A \land !B, \Gamma \fCenter \Delta$
\DisplayProof
\qquad
\Axiom$!B, \Gamma \fCenter \Delta$
\RightLabel{$\LeftR{\land}_2$}
\UnaryInf$!A \land !B, \Gamma \fCenter \Delta$
\DisplayProof
\]
Dene \Log{G3c} mung duwe siji aturan:
\[
\Axiom$ !A, !B, \Gamma \fCenter \Delta$
\RightLabel{$\LeftR{\land}_3$}
\UnaryInf$ !A \land !B, \Gamma \fCenter \Delta$
\DisplayProof
\]
Saiki kita bisa takon, ``Apa aturan saka \Log{G3c}, yaiku
$\LeftR{\land}_3$, admisibel ing \Log{G1c}?'' Wangsulane iya: saben inferensi nganggo
$\LeftR{\land}_3$ bisa disimulasi langsung mung nganggo aturan~\Log{G1c}.
Saliyane $\LeftR{\land}_1$ lan $\LeftR{\land}_2$, kita uga mbutuhake
aturan struktural~$\LeftR{\Contraction}$:
\[
\Axiom$!A, !B, \Gamma \fCenter \Delta$
\RightLabel{$\LeftR{\land}_1$}
\UnaryInf$!A \land !B, !B, \Gamma \fCenter \Delta$
\RightLabel{$\LeftR{\land}_2$}
\UnaryInf$!A \land !B, !A \land !B, \Gamma \fCenter \Delta$
\RightLabel{\LeftR{\Contraction}}
\UnaryInf$!A \land !B, \Gamma \fCenter \Delta$
\DisplayProof
\]
\end{ex}

Nyimulasi inferensi saka siji aturan nganggo aturan liyane minangka siji
cara nuduhake yen aturan kasebut admisibel. Nanging ana cara liyane,
lan sawetara aturan admisibel ora bisa disimulasi kaya mangkono.
Aturan sing bisa disimulasi diarani aturan sing \emph{bisa diturunake}.

\begin{defn}
Aturan~$R$,
\[
\AxiomC{$S_1$}
\AxiomC{$\dots$}
\AxiomC{$S_n$}
\RightLabel{$R$}
\TrinaryInfC{$S$}
\DisplayProof
\]
diarani \emph{bisa diturunake} ing sistem yen ana !!{proof} skematis
kanggo sekuen~$S$ nganggo aturan lan aksioma sistem kasebut bebarengan
karo sekuen awal tambahan sing wujude~$S_1$, \dots,~$S_n$.
\end{defn}

!!{proof} ing \olref{ex:land-deriv} nuduhake yen
$\LeftR{\land}_3$ minangka aturan sing bisa diturunake ing~\Log{G1c},
amarga !!{proof} kasebut minangka !!{proof} skematis kanggo kesimpulan
$\LeftR{\land}_3$ saka premis-premis $\LeftR{\land}_3$ minangka sekuen
awal, mung nganggo aturan \Log{G1c}. (Ora ana aksioma \Log{G1c} sing
digunakake.)

\begin{prop}\ollabel{prop:lor-G3-adm}
Aturan $\LeftR{\land}$ lan $\RightR{\lor}$ saka~\Log{G3c} bisa
diturunake ing~\Log{G1c}.
\end{prop}

\begin{proof}
  Bukti kanggo $\LeftR{\land}$ ana ing \olref{ex:land-deriv}.
  Bukti kanggo $\RightR{\lor}$ ditinggal minangka latihan.
\end{proof}

\begin{prob}
Tuduhake yen aturan $\RightR{\lor}$ saka \Log{G3c} bisa diturunake
ing~\Log{G1c}.
\end{prob}

\begin{prob}
Upamane $\liff$ minangka !!{operator} sing didhefinisikake:
$!A \liff !B$ minangka cekakan kanggo $(!A \lif !B) \land (!B \lif !A)$.
Tuduhake yen aturan
\begin{enumerate}
\item \Axiom$\Gamma, !A, !B \fCenter \Delta$
\Axiom$\Gamma \fCenter \Delta, !A, !B$
\RightLabel{\LeftR{\liff}}
\BinaryInf$!A \liff !B, \Gamma \fCenter \Delta$
\DisplayProof
\item
\Axiom$!A, \Gamma \fCenter \Delta, !B$
\Axiom$!B, \Gamma \fCenter \Delta, !A$
\RightLabel{\RightR{\liff}}
\BinaryInf$\Gamma \fCenter \Delta, !A \liff !B$
\DisplayProof
\end{enumerate}
bisa diturunake ing (a)~\Log{G1c} lan (b)~\Log{G3c}.
\end{prob}

\begin{ex}\ollabel{ex:weak-adm}
Aturan $\RightR{\Weakening}$ saka \Log{G1c},
\[
\Axiom$ \Gamma \fCenter \Delta$
\RightLabel{\RightR{\Weakening}}
\UnaryInf$ \Gamma \fCenter \Delta, !A$
\DisplayProof
\]
admisibel ing~\Log{G3c}. Gagasan buktine prasaja: upamane ana
!!a{proof}~$\pi$ kanggo premis $\Gamma \Sequent \Delta$ ing~\Log{G3c}.
Dhisik, ganti jeneng eigenparameter ing saben subwit bukti sing
gegayutan supaya seger tumrap formula sing bakal ditambahake; cara iki
njaga sekuen pungkasan lan dhuwure bukti. Tambahake !!{formula}~$!A$
ing sukseden saben sekuen ing~$\pi$. Aksioma tetep dadi aksioma;
inferensi sing bener tetep bener. !!{formula} pungkasan saka !!{proof}
anyar yaiku $\Gamma \Sequent \Delta, !A$.

Nanging aturan $\RightR{\Weakening}$ ora bisa diturunake ing~\Log{G3c}.
Upamane aturan iki bisa diturunake, yaiku upamane bisa ditemokake
!!a{proof} kanggo $\Gamma \Sequent \Delta, !A$ mung nganggo aksioma lan
aturan~\Log{G3c}, bisa uga kanthi sekuen awal tambahan~$\Gamma \Sequent
\Delta$. Yen !!{proof} iki murni skematis, kaya sing disaratake supaya
$\RightR{\Weakening}$ bisa diturunake, kita bisa ngowahi dadi !!a{proof}
kanggo~$\Sequent !A$ kanthi mbusak $\Gamma$ lan~$\Delta$. Iki ngowahi
sekuen awal tambahan~$\Gamma \Sequent \Delta$ dadi sekuen kosong~$\quad
\Sequent \quad$. Amarga sekuen kosong ora ngemot !!{formula}s babar
pisan, dene saben aturan~\Log{G3c} mbutuhake paling ora siji !!{formula}
aktif ing premis-premise, sekuen kosong ora bisa dadi premis inferensi
apa wae ing !!{proof} skematis iki. Mula !!{proof} kasebut mung bisa
skematis yen ora nggunakake sekuen tambahan $\Gamma \Sequent \Delta$
minangka sekuen awal. Kanthi tembung liya, iki bakal dadi derivasi
skematis kanggo $\Gamma \Sequent \Delta, !A$ mung nganggo aksioma lan
aturan~\Log{G3c}, lan kita bakal wis nuduhake yen saben sekuen awujud
iki duwe !!a{proof} ing~\Log{G3c}. Nanging iki mokal, amarga \Log{G3c}
sahih lan mung bisa !!{prove} sekuen sing bener ing saben !!{structure}.
Yen $\Gamma = \Delta = \emptyset$ lan $!A \ident \lfalse$, \Log{G3c}
kudu bisa !!{prove}, upamane, sekuen $\quad \Sequent \lfalse$, kamangka
sekuen iki ora bisa dibuktekake ing sistem kasebut.
\end{ex}

Saiki kita menehi bukti induktif yen pelemahan minangka aturan
admisibel ing~\Log{G3c}. Malah ana asil sing rada luwih kuwat, kaya
sing dituduhake argumentasi intuisi mau: kanthi nambahake $!A$ ing
sisih tengen saben sekuen ing~$\pi$, kita ora ngowahi struktur
!!{proof}; mligine, dhuwure ora mundhak. Amarga iki wigati, kita
menehi dhefinisi dhewe.

\begin{defn}
Aturan~$R$,
\[
\AxiomC{$S_1$}
\AxiomC{$\dots$}
\AxiomC{$S_n$}
\RightLabel{$R$}
\TrinaryInfC{$S$}
\DisplayProof
\]
diarani \emph{admisibel kanthi njaga !!{height}} (utawa
\emph{admisibel-!!{hp}}) ing sistem yen, saben $\Proves[h_i] S_i$
kanggo $i = 1$, \dots,~$n$, uga $\Proves_m S$ kanthi
$m = \max(h_1, \dots, h_n)$.
\end{defn}

Gathekna yen supaya aturan admisibel kanthi njaga !!{height}, cukup
yen saben premis~$S_i$ duwe !!{proof}s~$\pi_i$ kanthi
!!{height}~$h_i = \pheight{\pi_i}$, ana !!a{proof}~$\pi$ kanggo
kesimpulan kanthi !!{height}~$\pheight{\pi}$ sing ora ngluwihi nilai
paling gedhe saka~$h_i$. Mligine, yen mung ana siji premis~$S_1$,
cukup yen $\pheight{\pi} \le \pheight{\pi_1}$. Kanggo
$\RightR{\Weakening}$ lan $\LeftR{\Weakening}$, malah
$\pheight{\pi} = \pheight{\pi_1}$, nanging pepadhan iki ora wajib.

\begin{prop}\ollabel{prop:weak-G3c-adm}
Aturan $\RightR{\Weakening}$ lan $\LeftR{\Weakening}$ saka \Log{G1c}
admisibel kanthi njaga !!{height} ing~\Log{G3c}.
\end{prop}

\begin{proof}
  Kita kudu nuduhake yen $\Log{G3c} \Proves \Gamma \Sequent \Delta$
  kanthi !!{proof}~$\pi$ sing duwe !!{height}~$\pheight{\pi} = n$,
  mula ana \Log{G3c}-!!{proof} $\pi'$ kanggo $\Gamma \Sequent \Delta,
  !A$ kanthi $\pheight{\pi'} = n$. Dhisik, pilih eigenparameter sing
  seger tumrap formula tambahan kanthi ngganti jeneng ing subwit
  kaya sing wis dijlentrehake. Kita mbuktekake kanthi induksi
  tumrap~$n$. Yen $n=0$, $\pi$ mung aksioma $\Gamma \Sequent \Delta$
  (yaiku ana formula atomik $!C$ ing $\Gamma$ lan~$\Delta$, utawa ana
  falsitas ing anteseden), lan $\Gamma \Sequent \Delta, !A$ uga
  aksioma. Yen $n = \pheight{\pi} >0$, !!{proof}~$\pi$ duwe inferensi
  pungkasan. Kita misahake kasus miturut aturan sing digunakake.
  Kita mung menehi siji conto: upamane inferensi pungkasan nggunakake
  $\RightR{\land}$. Mula $\Delta = \Delta', !B \land !C$, lan
  sub-!!{proof}s langsung $\pi_1$ lan~$\pi_2$ saka inferensi pungkasan
  duwe sekuen pungkasan $\Gamma \Sequent \Delta', !B$ lan
  $\Gamma \Sequent \Delta', !C$, manut urutane. Kanthi tembung liya,
  $\pi$ duwe wujud iki:
  \[
  \AxiomC{}
  \RightLabel{$\pi_1$}
  \Deduce$\Gamma \fCenter \Delta', !B$
  \AxiomC{}
  \RightLabel{$\pi_2$}
  \Deduce$\Gamma \fCenter \Delta', !C$
  \RightLabel{\RightR{\land}}
  \BinaryInf$\Gamma \fCenter \Delta', !B \land !C$
  \DisplayProof
  \]
  Kita uga duwe $n = \max(\pheight{\pi_1}, \pheight{\pi_2}) + 1$.
  Dadi $\pheight{\pi_1} < n$ lan $\pheight{\pi_2} < n$, lan
  hipotesis induksi bisa ditrapake: ana !!{proof}s $\pi_1'$ lan~$\pi_2'$
  kanggo $\Gamma \Sequent \Delta', !B, !A$ lan $\Gamma \Sequent
  \Delta', !C, !A$. Kanthi aturan $\RightR{\land}$, kita oleh
  !!a{proof}~$\pi'$ kanggo $\Gamma \Sequent \Delta', !B \land !C, !A$:
  \[
  \AxiomC{}
  \RightLabel{$\pi_1'$}
  \Deduce$\Gamma \fCenter \Delta', !B, !A$
  \AxiomC{}
  \RightLabel{$\pi_2'$}
  \Deduce$\Gamma \fCenter \Delta', !C, !A$
  \RightLabel{\RightR{\land}}
  \BinaryInf$\Gamma \fCenter \Delta', !B \land !C, !A$
  \DisplayProof
  \]
  Kita duwe $\pheight{\pi'} = \max(\pheight{\pi_1'}, \pheight{\pi_2'})
  + 1 = \max(\pheight{\pi_1}, \pheight{\pi_2}) + 1 = \pheight{\pi}$.
\end{proof}

\olref{prop:weak-G3c-adm} migunani banget. Kita ngenalake notasi
kanggo ngetrapake asil iki bola-bali: yen $\pi$ minangka !!a{proof}
kanggo~$\Gamma \Sequent \Delta$, mula $\pi[\Pi \Sequent \Lambda]$
minangka asil ngetrapake admisibilitas pelemahan kiwa kanggo kabeh
!!{formula}s ing~$\Pi$ lan admisibilitas pelemahan tengen kanggo kabeh
!!{formula}s ing~$\Lambda$. Iki minangka !!a{proof} kanggo
$\Gamma, \Pi \Sequent \Delta, \Lambda$.

\end{document}
